% GENERATED FILE. Edit canonical lessons, then run npm run generate:books.
% canonical-source-hash: fnv1a64:5970b9ac45aae1a2
% canonical-lessons: GE-C280-verbrechen, GE-C280-verbrecher, GE-C280-taeter, GE-C280-opfer, GE-C280-tatort

\chapter{Das Verbrechen}
\label{ch:ge-das-verbrechen}
\hlchaptermodality{280}

\begin{chapteropening}
\textbf{By the end of this chapter:} \emph{I can talk about a crime: a crime, a criminal, an offender, a victim and the crime scene.}

Complete the last lesson of chapter 280: I can talk about a crime: a crime, a criminal, an offender, a victim and the crime scene.
\end{chapteropening}

\section[das Verbrechen]{das Verbrechen — the crime}

\label{lesson:GE-C280-verbrechen}



\begin{quote}
Before the new one: say the German for to happen, then the German for the incident.
\end{quote}

\subsection*{You'll want to know: das Verbrechen}
\textbf{das Verbrechen} — \textquotedblleft{}the crime\textquotedblright{}.

\begin{quote}
Das Verbrechen geschah in der Nacht.
\end{quote}

\begin{quote}
\emph{The crime happened at night.}
\end{quote}

Plural: \emph{die Verbrechen}. For a serious crime; a lesser offence is \emph{ein Vergehen}.

\begin{cousinweb}
\emph{Verbrechen} is the old verb \emph{verbrechen} \textquotedblleft{}to do wrong\textquotedblright{} used as a noun, built on \emph{brechen} \textquotedblleft{}to break\textquotedblright{}, Proto-Germanic \emph{brekaną} (EtymDB): to break the law.
\end{cousinweb}

\begin{grammarlens}[title={What you've built}]
One more word for telling, explaining and describing at B1.
\end{grammarlens}

\subsection*{Your turn}
\begin{itemize}
\raggedright
  \item \emph{Say it:} \emph{das Verbrechen}
  \item \emph{Say it:} \emph{das Verbrechen}, once more
  \item \emph{Say it:} a sentence of your own with \emph{das Verbrechen}
  \item \emph{Recall:} say the German for to happen, then the German for the incident, then say \emph{das Verbrechen} again
\end{itemize}

\begin{tcolorbox}[breakable,colback=teal!4,colframe=teal!35!black,title={Before you move on}]
What does das Verbrechen mean? (\textquotedblleft{}The crime\textquotedblright{}.) Say it once more.
\end{tcolorbox}
\section[der Verbrecher]{der Verbrecher — the criminal}

\label{lesson:GE-C280-verbrecher}



\begin{quote}
Before the new one: say the German for the incident, then the German for the crime.
\end{quote}

\subsection*{You'll want to know: der Verbrecher}
\textbf{der Verbrecher} — \textquotedblleft{}the criminal\textquotedblright{}.

\begin{quote}
Die Polizei suchte den Verbrecher drei Wochen lang.
\end{quote}

\begin{quote}
\emph{The police searched for the criminal for three weeks.}
\end{quote}

Plural: \emph{die Verbrecher}. Feminine \emph{die Verbrecherin}.

\begin{cousinweb}
\emph{Verbrecher} is built on \emph{verbrechen} \textquotedblleft{}to do wrong\textquotedblright{}, from \emph{brechen} \textquotedblleft{}to break\textquotedblright{}, Proto-Germanic \emph{brekaną} (EtymDB), with the ending \emph{-er} for a person.
\end{cousinweb}

\begin{grammarlens}[title={What you've built}]
One more word for telling, explaining and describing at B1.
\end{grammarlens}

\subsection*{Your turn}
\begin{itemize}
\raggedright
  \item \emph{Say it:} \emph{der Verbrecher}
  \item \emph{Say it:} \emph{der Verbrecher}, once more
  \item \emph{Say it:} a sentence of your own with \emph{der Verbrecher}
  \item \emph{Recall:} say the German for the incident, then the German for the crime, then say \emph{der Verbrecher} again
\end{itemize}

\begin{tcolorbox}[breakable,colback=teal!4,colframe=teal!35!black,title={Before you move on}]
What does der Verbrecher mean? (\textquotedblleft{}The criminal\textquotedblright{}.) Say it once more.
\end{tcolorbox}
\section[der Täter]{der Täter — the offender, the culprit}

\label{lesson:GE-C280-taeter}



\begin{quote}
Before the new one: say the German for the crime, then the German for the criminal.
\end{quote}

\subsection*{You'll want to know: der Täter}
\textbf{der Täter} — \textquotedblleft{}the offender, the culprit\textquotedblright{}.

\begin{quote}
Der Täter trug eine schwarze Mütze.
\end{quote}

\begin{quote}
\emph{The offender wore a black cap.}
\end{quote}

Plural: \emph{die Täter}. Feminine \emph{die Täterin}. The person who did the deed, \emph{die Tat}.

\begin{cousinweb}
\emph{Täter} is built on \emph{Tat} \textquotedblleft{}deed\textquotedblright{}, from Old High German \emph{tāt}, Proto-Germanic \emph{dēdiz} (EtymDB), a cousin of English \emph{deed}.
\end{cousinweb}

\begin{grammarlens}[title={What you've built}]
One more word for telling, explaining and describing at B1.
\end{grammarlens}

\subsection*{Your turn}
\begin{itemize}
\raggedright
  \item \emph{Say it:} \emph{der Täter}
  \item \emph{Say it:} \emph{der Täter}, once more
  \item \emph{Say it:} a sentence of your own with \emph{der Täter}
  \item \emph{Recall:} say the German for the crime, then the German for the criminal, then say \emph{der Täter} again
\end{itemize}

\begin{tcolorbox}[breakable,colback=teal!4,colframe=teal!35!black,title={Before you move on}]
What does der Täter mean? (\textquotedblleft{}The offender, the culprit\textquotedblright{}.) Say it once more.
\end{tcolorbox}
\section[das Opfer]{das Opfer — the victim; the sacrifice}

\label{lesson:GE-C280-opfer}



\begin{quote}
Before the new one: say the German for the criminal, then the German for the offender.
\end{quote}

\subsection*{You'll want to know: das Opfer}
\textbf{das Opfer} — \textquotedblleft{}the victim; the sacrifice\textquotedblright{}.

\begin{quote}
Das Opfer konnte den Mann gut beschreiben.
\end{quote}

\begin{quote}
\emph{The victim could describe the man well.}
\end{quote}

Plural: \emph{die Opfer}. Neuter for anyone: \emph{Er war das Opfer}.

\begin{cousinweb}
\emph{Opfer} is inherited from Old High German \emph{opfar} \textquotedblleft{}offering\textquotedblright{} (EtymDB), which is usually traced to Church Latin \emph{operārī} \textquotedblleft{}to work, to make an offering\textquotedblright{}.
\end{cousinweb}

\begin{grammarlens}[title={What you've built}]
One more word for telling, explaining and describing at B1.
\end{grammarlens}

\subsection*{Your turn}
\begin{itemize}
\raggedright
  \item \emph{Say it:} \emph{das Opfer}
  \item \emph{Say it:} \emph{das Opfer}, once more
  \item \emph{Say it:} a sentence of your own with \emph{das Opfer}
  \item \emph{Recall:} say the German for the criminal, then the German for the offender, then say \emph{das Opfer} again
\end{itemize}

\begin{tcolorbox}[breakable,colback=teal!4,colframe=teal!35!black,title={Before you move on}]
What does das Opfer mean? (\textquotedblleft{}The victim; the sacrifice\textquotedblright{}.) Say it once more.
\end{tcolorbox}
\section[der Tatort]{der Tatort — the crime scene}

\label{lesson:GE-C280-tatort}



\begin{quote}
Before the new one: say the German for the offender, then the German for the victim.
\end{quote}

\subsection*{You'll want to know: der Tatort}
\textbf{der Tatort} — \textquotedblleft{}the crime scene\textquotedblright{}.

\begin{quote}
Die Polizei kam schnell zum Tatort.
\end{quote}

\begin{quote}
\emph{The police came quickly to the crime scene.}
\end{quote}

Plural: \emph{die Tatorte}. \emph{Tat} 'deed' plus \emph{Ort} 'place'.

\begin{cousinweb}
\emph{Tatort} joins \emph{Tat} \textquotedblleft{}deed\textquotedblright{}, Proto-Germanic \emph{dēdiz}, and \emph{Ort} \textquotedblleft{}place\textquotedblright{}, Old High German \emph{ort} \textquotedblleft{}point, place\textquotedblright{} (EtymDB): the place of the deed.
\end{cousinweb}

\begin{grammarlens}[title={What you've built}]
One more word for telling, explaining and describing at B1.
\end{grammarlens}

\subsection*{Your turn}
\begin{itemize}
\raggedright
  \item \emph{Say it:} \emph{der Tatort}
  \item \emph{Say it:} \emph{der Tatort}, once more
  \item \emph{Say it:} a sentence of your own with \emph{der Tatort}
  \item \emph{Recall:} say the German for the offender, then the German for the victim, then say \emph{der Tatort} again
\end{itemize}

\begin{tcolorbox}[breakable,colback=teal!4,colframe=teal!35!black,title={Before you move on}]
What does der Tatort mean? (\textquotedblleft{}The crime scene\textquotedblright{}.) Say it once more.
\end{tcolorbox}
